\documentclass[10pt,a4paper]{article}
\usepackage{lmodern}
\usepackage[margin=19mm]{geometry}
\usepackage{booktabs}
\usepackage{graphicx}
\usepackage[table]{xcolor}
\usepackage[hidelinks]{hyperref}
\usepackage{parskip}
\usepackage{array}
\graphicspath{{figures/}}
\definecolor{warning}{HTML}{FFF3CD}
\definecolor{info}{HTML}{E0F2FE}
\hypersetup{pdftitle={Dynamic and static comparison of three UWB protocols},pdfauthor={UWB ESP-IDF technical report}}
\begin{document}
\begin{titlepage}
\centering
\vspace*{18mm}
{\Huge\bfseries Dynamic and static UWB comparison\par}
\vspace{4mm}
{\Large FlexTDOA, Native DS-TWR and Passive DS-TWR\par}
\vspace{14mm}
{\large Field campaign of August 6, 2026\par}
\vspace{12mm}
\begin{tabular}{rl}
Platform: & 5 ESP32-S3 + DW3000 modules, Raspberry Pi collector\\
Topology: & M1 tag; M2, M3, M4 and M5 anchors\\
Radio: & UWB channel 9; configured 40 MHz SPI\\
Firmware branch: & \texttt{imu-500hz-fusion}\\
Evidence: & six complete immutable JSONL captures plus replay summaries\\
\end{tabular}
\vfill
\colorbox{info}{\parbox{0.88\textwidth}{\vspace{2mm}
\textbf{Main result.} All three protocols produced usable unfiltered dynamic
trajectories. FlexTDOA delivered the highest independent dynamic rate and the
best static local precision. Native DS-TWR had the shortest worst-case dynamic
gap. Passive DS-TWR remained usable, but its overlapping publications must be
reduced to independent frames for a fair rate comparison.\vspace{2mm}}}
\vfill
{\large Reproducible technical report\par}
\end{titlepage}
\pagenumbering{roman}
\tableofcontents
\clearpage
\pagenumbering{arabic}

\section{Executive summary}
\colorbox{warning}{\parbox{0.96\textwidth}{\textbf{Evidence boundary.}
Stream rates, uptime gaps, sample counts, fusion acceptance and reset counts are
direct diagnostics. Absolute UWB--RTK disagreement is only indicative: RTK was
not an independent survey, clocks are associated in software, antenna lever
arms are unknown, and the UWB frame is registered through RTK anchor geometry.}}

Passive DS overlapping raw windows remain visible in the telemetry rate, while
all comparative precision and trajectory metrics use independent frames.
Interleaved records marked \texttt{imu\_fused=true} are excluded from raw UWB
metrics. Outliers are retained.

The dynamic independent rates were \textbf{27.62 positions/s} for FlexTDOA, \textbf{22.38 positions/s} for Native DS-TWR, and \textbf{23.84 positions/s} for Passive DS-TWR. Passive DS emitted 71.33 dashboard events/s, but those include overlapping windows.

The static raw local RMS values were \textbf{2.13 cm} (FlexTDOA), \textbf{2.02 cm} (Native DS-TWR) and \textbf{1.78 cm} (Passive DS-TWR). These are precision around each capture median, not absolute accuracy.


The stationary tag RTK centers differed by 1.583 m between blocks, invalidating
an absolute cross-protocol static ranking. IMU transport was healthy at roughly
119--127 telemetry samples/s, but all six replays reported zero yaw-alignment
updates. This campaign therefore does not demonstrate acceleration-derived
position improvement.

\section{Scope and questions}
The campaign asks how many independent positions each protocol delivers, where
stream gaps occur, how stable the unfiltered stationary position is, and whether
the first IMU-assisted replay improves the result without hiding raw UWB.

The July 26 report used known static geometry and could make a direct accuracy
statement. Here the anchors and tag are tied to mobile RTK-derived geometry and
software time association, so the evidence boundary is deliberately stricter.

\section{System and protocol background}
\subsection{FlexTDOA}
The tag is UWB receive-only. Anchor transmissions form directed time-difference
observations and the tag solves from a recent coherent observation set. Radio
airtime is largely independent of the number of passive tags.

\subsection{Native DS-TWR}
The active tag exchanges poll, response and final messages with every anchor,
obtains four absolute ranges and solves a coherent frame. Airtime grows with
active tag count, but frame semantics are direct.

\subsection{Passive DS-TWR}
The tag listens to anchor-to-anchor DS-TWR exchanges. Rolling windows may
produce several publications from shared measurements, so dashboard events and
independent frames are intentionally reported separately.

\section{Test configuration and captured evidence}
The dynamic walks lasted 91--96 s and the stationary blocks approximately 180 s.
The selector uses the final walk and the RTK-fixed Raspberry static capture for
each protocol.
\begin{center}\scriptsize\resizebox{\textwidth}{!}{%
\begin{tabular}{lrrrrrrrr}\toprule
\textbf{Protocol} & \textbf{Mode} & \textbf{Capture ID} & \textbf{duration s} & \textbf{events} & \textbf{independent} & \textbf{indep Hz} & \textbf{coverage pct} & \textbf{complete} \\
\midrule
FlexTDOA & dynamic & flextdoa\_walk\_imu\_clean & 150.2 & 4149 & 4149 & 27.62 & 99.9 & yes \\
Native DS-TWR & dynamic & native\_ds\_walk\_imu\_clean\_take2 & 150.1 & 3359 & 3359 & 22.38 & 99.9 & yes \\
Passive DS-TWR & dynamic & passive\_ds\_walk\_imu\_clean & 150.0 & 10702 & 3576 & 23.84 & 99.9 & yes \\
FlexTDOA & static & flextdoa\_static\_final & 30.3 & 929 & 929 & 30.67 & 99.4 & yes \\
Native DS-TWR & static & native\_ds\_static\_final & 30.5 & 686 & 686 & 22.47 & 99.4 & yes \\
Passive DS-TWR & static & passive\_ds\_static\_final & 31.1 & 2295 & 772 & 24.80 & 99.1 & yes \\
\bottomrule\end{tabular}}\end{center}
Every selected block has a terminal \texttt{capture\_end}. All raw data are
archived losslessly and identified by SHA-256.

\section{Methodology}
\subsection{Population and continuity}
No error-based trimming or temporal smoothing is applied. Rates use full
capture wall time. Gaps are positive differences between unique ESP32 uptime
values; P50/P95/P99, maximum and counts above 100 ms are retained together.

\subsection{Static precision}
Radial displacement is computed around each block's component-wise position
median. CEP50, RMS, P95, P99, maximum, 2DRMS and first-to-last-decile drift use
all independent samples. They describe repeatability, not offset from truth.

\subsection{RTK audit}
Only quality-4 fixes enter the audit. WGS84 is converted through ECEF to a common
ENU frame. A proper 2-D rotation and translation is fitted from at least three
anchor centers; scale is fixed. Tag matches must be within 100 ms. Anchor-fit
residual and association timing accompany every tag error.

\section{Dynamic results}
\subsection{Raw trajectories}
The walks are separate physical trials, not synchronized repeats. The panels
therefore show continuity and gross geometry rather than point-by-point route
agreement.

\begin{figure}[htbp]
\centering
\includegraphics[width=\textwidth]{figures/06_dynamic_trajectories.pdf}
\caption{Raw independent dynamic trajectories in each local UWB frame. Green markers are anchors. Different panels represent different walks.}
\end{figure}
\subsection{Update rate and gaps}
\begin{center}\scriptsize\resizebox{\textwidth}{!}{%
\begin{tabular}{lrrrrrrrr}\toprule
\textbf{Protocol} & \textbf{duration s} & \textbf{events} & \textbf{independent} & \textbf{event Hz} & \textbf{indep Hz} & \textbf{gap P95 ms} & \textbf{gap max ms} & \textbf{step P95 m/s} \\
\midrule
FlexTDOA & 150.2 & 4149 & 4149 & 27.62 & 27.62 & 70 & 2030 & 5.08 \\
Native DS-TWR & 150.1 & 3359 & 3359 & 22.38 & 22.38 & 60 & 410 & 5.20 \\
Passive DS-TWR & 150.0 & 10702 & 3576 & 71.33 & 23.84 & 30 & 610 & 4.60 \\
\bottomrule\end{tabular}}\end{center}
\begin{figure}[htbp]
\centering
\includegraphics[width=0.94\textwidth]{figures/01_position_rate.pdf}
\caption{Independent position rate in the selected dynamic and static blocks. Passive DS all-event publications are not substituted for independent frames.}
\end{figure}

\begin{figure}[htbp]
\centering
\includegraphics[width=0.94\textwidth]{figures/02_position_gap.pdf}
\caption{P95 positive gap between unique position uptimes. Maximum gaps and counts remain in the tables.}
\end{figure}

Native DS-TWR had the best dynamic worst-case continuity (200 ms maximum).
FlexTDOA produced the highest independent rate but one 520 ms gap. Passive DS
had a useful typical cadence yet one 1.61 s interruption and the most gaps over
100 ms. Its recovery is operationally useful, but rare long gaps remain the
first diagnostic target. Step speed is a sample-to-sample trajectory diagnostic,
not measured walking speed.

\subsection{Complete continuity audit}
\begin{center}\scriptsize\resizebox{\textwidth}{!}{%
\begin{tabular}{lrrrrrrrr}\toprule
\textbf{Protocol} & \textbf{Mode} & \textbf{event Hz} & \textbf{indep Hz} & \textbf{P50 ms} & \textbf{P95} & \textbf{P99} & \textbf{max} & \textbf{>100 ms} \\
\midrule
FlexTDOA & dynamic & 27.62 & 27.62 & 30 & 70 & 120 & 2030 & 57 \\
Native DS-TWR & dynamic & 22.38 & 22.38 & 40 & 60 & 70 & 410 & 3 \\
Passive DS-TWR & dynamic & 71.33 & 23.84 & 10 & 30 & 130 & 610 & 172 \\
FlexTDOA & static & 30.67 & 30.67 & 30 & 60 & 88 & 140 & 1 \\
Native DS-TWR & static & 22.47 & 22.47 & 40 & 50 & 60 & 70 & 0 \\
Passive DS-TWR & static & 73.72 & 24.80 & 10 & 20 & 138 & 190 & 38 \\
\bottomrule\end{tabular}}\end{center}

\section{Static results}
\subsection{Raw local precision}
\begin{center}\scriptsize\resizebox{\textwidth}{!}{%
\begin{tabular}{lrrrrrrrr}\toprule
\textbf{Protocol} & \textbf{N} & \textbf{CEP50 cm} & \textbf{RMS cm} & \textbf{P95 cm} & \textbf{P99 cm} & \textbf{max cm} & \textbf{2DRMS cm} & \textbf{drift cm} \\
\midrule
FlexTDOA & 929 & 1.68 & 2.13 & 3.85 & 4.89 & 5.95 & 4.27 & 0.21 \\
Native DS-TWR & 686 & 1.40 & 2.02 & 3.76 & 5.99 & 10.06 & 4.02 & 0.16 \\
Passive DS-TWR & 772 & 1.50 & 1.78 & 3.11 & 3.72 & 6.44 & 3.57 & 0.45 \\
\bottomrule\end{tabular}}\end{center}
\begin{figure}[htbp]
\centering
\includegraphics[width=0.94\textwidth]{figures/03_static_precision.pdf}
\caption{Static precision around each block's own median. All independent raw samples are included.}
\end{figure}

\begin{figure}[htbp]
\centering
\includegraphics[width=\textwidth]{figures/07_static_position_clouds.pdf}
\caption{Static raw position clouds centered on their own medians and drawn at one common scale. The cross is the median.}
\end{figure}

\begin{figure}[htbp]
\centering
\includegraphics[width=\textwidth]{figures/08_static_gps_rtk_clouds.pdf}
\caption{GPS RTK Fixed tag solutions. Top: each protocol block is centered on its own median with one shared centimetre scale, showing within-block precision. Bottom: the same fixes and medians in the common ENU frame with equal axis scale, exposing the 1.58 m cross-capture reference jump.}
\end{figure}

FlexTDOA has the smallest static RMS and P95. Its isolated 55.12 cm maximum is
retained rather than trimmed. Native DS has the broadest P95, while Passive DS
lies between them on RMS and has the smallest maximum. One stationary point
cannot characterize GDOP or spatial bias across the anchor polygon.

The RTK Fixed clouds are individually tight: their within-block P95 horizontal
spreads are approximately 0.85 cm, 1.15 cm and 2.02 cm for FlexTDOA, Native DS
and Passive DS respectively. The common-ENU view nevertheless separates the
Passive DS block by about 1.58 m. High short-term RTK precision therefore does
not imply cross-session correctness.

\subsection{RTK cross-capture consistency gate}
The stationary RTK center's maximum pairwise separation is \textbf{2.032 m}, versus a 0.250 m gate. \textbf{static absolute UWB-RTK errors are audit-only and invalid for cross-protocol ranking}. FlexTDOA and Native DS centers differ by about 6.5 mm, whereas the Passive DS block differs by approximately 1.58 m.


\section{IMU capture and fusion assessment}
The sensor request was configured at 200 Hz; the collector observed roughly
119--127 events/s. This is a telemetry result and does not prove that every
hardware sample reached the Raspberry Pi.
\begin{center}\scriptsize\resizebox{\textwidth}{!}{%
\begin{tabular}{lrrrrrr}\toprule
\textbf{Protocol} & \textbf{Mode} & \textbf{IMU events} & \textbf{observed Hz} & \textbf{valid pct} & \textbf{gap P95 ms} & \textbf{gap max ms} \\
\midrule
FlexTDOA & dynamic & 14979 & 99.7 & 100.0 & 20 & 60 \\
Native DS-TWR & dynamic & 14937 & 99.5 & 100.0 & 20 & 50 \\
Passive DS-TWR & dynamic & 14993 & 99.9 & 100.0 & 30 & 80 \\
FlexTDOA & static & 13330 & 440.1 & 99.0 & 10 & 30 \\
Native DS-TWR & static & 13586 & 445.0 & 99.2 & 10 & 40 \\
Passive DS-TWR & static & 13652 & 438.5 & 99.1 & 20 & 60 \\
\bottomrule\end{tabular}}\end{center}\subsection{Replay result}
Replay acceptance, outliers and resets are direct diagnostics. RTK columns are
secondary evidence because repeated output matches do not create independent
fixes.
\begin{center}\scriptsize\resizebox{\textwidth}{!}{%
\begin{tabular}{lrrrrrrrrr}\toprule
\textbf{Protocol} & \textbf{Mode} & \textbf{accept pct} & \textbf{outliers} & \textbf{resets} & \textbf{raw RTK m} & \textbf{fused RTK m} & \textbf{delta mm} & \textbf{yaw updates} & \textbf{anchor fit m} \\
\midrule
FlexTDOA & dynamic & 99.8 & 14 & 0 & 0.435 & 0.435 & -0.4 & 4 & 0.009 \\
Native DS-TWR & dynamic & 96.4 & 165 & 0 & 0.478 & 0.478 & -0.1 & 1 & 0.020 \\
Passive DS-TWR & dynamic & 97.7 & 129 & 0 & 0.440 & 0.440 & -0.3 & 3 & 0.002 \\
FlexTDOA & static & 100.0 & 0 & 0 & n/a & n/a & n/a & 0 & n/a \\
Native DS-TWR & static & 99.7 & 2 & 0 & 2.972 & 2.974 & 2.5 & 0 & 1.064 \\
Passive DS-TWR & static & 100.0 & 0 & 0 & 2.369 & 2.369 & 0.0 & 0 & 0.715 \\
\bottomrule\end{tabular}}\end{center}
All replays have zero yaw-alignment updates, so body-frame horizontal
acceleration cannot safely be rotated into the map. Dynamic raw/fused RMSE
changes only by millimetres. The fused stream is therefore continuity-oriented
prediction/correction, not demonstrated inertial dead reckoning. Raw UWB stays
separate and is the basis of this protocol comparison.

\section{UWB--RTK association audit}
The following values are retained for transparency, not promoted to surveyed
accuracy. Debiased RMSE removes median tag residual and is a shape diagnostic.
\begin{center}\scriptsize\resizebox{\textwidth}{!}{%
\begin{tabular}{lrrrrrrrr}\toprule
\textbf{Protocol} & \textbf{Mode} & \textbf{pairs} & \textbf{time P95 ms} & \textbf{RMSE m} & \textbf{P95 m} & \textbf{debiased m} & \textbf{anchor fit m} & \textbf{alignment} \\
\midrule
FlexTDOA & dynamic & 98 & 36.1 & 0.437 & 0.830 & 0.445 & 0.013 & usable\_with\_limitations \\
Native DS-TWR & dynamic & 91 & 28.4 & 0.471 & 1.180 & 0.464 & 0.015 & usable\_with\_limitations \\
Passive DS-TWR & dynamic & 90 & 47.7 & 0.464 & 0.992 & 0.455 & 0.013 & usable\_with\_limitations \\
FlexTDOA & static & 0 & n/a & n/a & n/a & n/a & n/a & unavailable \\
Native DS-TWR & static & 129 & 23.0 & 2.979 & 3.002 & 0.025 & 1.073 & poor\_alignment\_do\_not\_claim\_accuracy \\
Passive DS-TWR & static & 244 & 42.9 & 1.909 & 1.933 & 0.020 & 0.741 & poor\_alignment\_do\_not\_claim\_accuracy \\
\bottomrule\end{tabular}}\end{center}
\begin{figure}[htbp]
\centering
\includegraphics[width=0.94\textwidth]{figures/04_dynamic_rtk_error_cdf.pdf}
\caption{Dynamic UWB--RTK disagreement after anchor-based rigid registration. This is not surveyed position error.}
\end{figure}

\begin{figure}[htbp]
\centering
\includegraphics[width=0.94\textwidth]{figures/05_alignment_quality.pdf}
\caption{Anchor-registration residual beside paired tag disagreement. Registration uncertainty limits tag interpretation.}
\end{figure}

Dynamic paths and anchor fits differ by capture. Passive DS dynamic geometry is
partly GNSS-derived, making its registration non-independent. The RTK audit is
still useful for gross frame errors and timing failures.

\section{Protocol interpretation and selection}
\subsection{FlexTDOA}
Best independent dynamic rate and static raw precision in this campaign, with
a receive-only scalable tag. The 520 ms dynamic maximum gap is the principal
caveat. The walks do not justify another protocol rewrite.

\subsection{Native DS-TWR}
Clearest coherent-frame semantics and best dynamic maximum gap. Static P95 is
wider here, so per-link bias/antenna-delay validation remains worthwhile. It is
the preferred deterministic baseline for few active tags.

\subsection{Passive DS-TWR}
Successfully follows motion without tag transmissions. Only one third of its
rolling publications are independent in these captures; the 1.61 s dynamic gap
and higher replay reset count keep recovery diagnostics as the priority.

\begin{center}\small
\begin{tabular}{p{0.20\textwidth}p{0.23\textwidth}p{0.23\textwidth}p{0.23\textwidth}}
\toprule
Criterion & FlexTDOA & Native DS-TWR & Passive DS-TWR\\
\midrule
Tag radio role & receive-only & active exchanges & receive-only\\
Dynamic independent rate & highest & middle & lowest after overlap removal\\
Dynamic worst gap & 520 ms & best: 200 ms & 1.61 s\\
Static raw precision & best RMS/P95 & broadest P95 & middle\\
Scaling with tag count & strongest & airtime grows & strong\\
Recommendation & primary scalable mode & coherent baseline & continue recovery tuning\\
\bottomrule
\end{tabular}
\end{center}

\section{Limitations}
\begin{itemize}
\item RTK is an in-system comparison reference, not independently surveyed truth.
\item Measurement time subtracts receiver-reported fix age from collector time; clocks are not hardware synchronized.
\item GNSS and UWB antenna phase centers have unknown lever arms.
\item Dynamic routes are not repeated surveyed paths; static/dynamic blocks are single trials.
\item One static point does not cover edge geometry, NLOS, orientation or spatial bias.
\item Collector rate is end-to-end telemetry, not direct DW3000 PHY throughput.
\end{itemize}

\section{Recommended next tests}
\begin{enumerate}
\item Repeat every protocol on one marked route with identical orientation and synchronized markers.
\item Interleave static blocks at the center, edges and outside the polygon.
\item Keep one continuous RTK session and log correction age/reset state across protocol changes.
\item Add hardware-correlated timestamps for latency claims.
\item Validate body-to-map yaw before enabling horizontal acceleration corrections.
\item For Passive DS, capture focused diagnostics around every gap above 250 ms.
\end{enumerate}

\section{Preliminary findings}
\begin{itemize}
\item FlexTDOA is the strongest overall scalable result in this campaign.
\item Native DS-TWR remains the clean coherent-frame and continuity baseline.
\item Passive DS-TWR is usable, but overlap-aware rate accounting and rare-gap recovery remain essential.
\item The RTK consistency failure prevents a static absolute-accuracy ranking; local precision alone is ranked.
\item IMU transport is healthy, but acceleration-assisted accuracy is not yet demonstrated.
\end{itemize}

\section{Reproduction and audit files}
The directory contains \texttt{analysis\_summary.json}, capture/replay/alignment
CSV files, SVG/PDF plots and this TeX source. The machine-readable summary
records whether optional lossless RAW archive generation was enabled; when it
is skipped, immutable source JSONL paths remain recorded for reproduction.

\begin{verbatim}
python tools/uwb_dynamic_static_report.py --input-dir INPUT --output-dir OUTPUT
\end{verbatim}

No raw capture is modified during report generation.
\end{document}
